% THETA3-Q10 contribution by /root/theta_032, 21 November 2026.
% GPT-6 Astra Pro; inherited reasoning-effort label not separately exposed.

\subsection{THETA3-Q10: determinant normalization and a finite certificate in every degree}\label{Q10}

\paragraph{The exact Fock determinant norm.}
Let $u_1,\ldots,u_n\in\C^n$, put
$\ell_u(z)=\sum_a u_a z_a$, and let $z_1,\ldots,z_n$ be independent
standard complex Gaussian vectors, each with probability density
$\pi^{-n}e^{-\|z\|^2}$ relative to real Lebesgue measure. Define
\[
 P_u(z_1,\ldots,z_n)=\det(\ell_{u_i}(z_j))_{1\leq i,j\leq n},
 \qquad G_u=(\langle u_i,u_j\rangle)_{i,j},
\]
with the Hermitian inner product conjugate linear in its first
variable. Then
\begin{equation}\label{Q10-Fock}
 \|P_u\|^2=n!\det G_u,\qquad
 \langle P_u,P_v\rangle
       =n!\det(\langle u_i,v_j\rangle)_{i,j}.
\end{equation}

\paragraph{Proof.}
The Gaussian covariance is
$\mathbb E[\overline{z_a}z_b]=\delta_{ab}$, so
$\mathbb E[\overline{\ell_u(z)}\ell_v(z)]
=\langle u,v\rangle$.
Expand both determinants. Independence of the $n$ columns gives
\[
 \langle P_u,P_v\rangle
 =\sum_{\sigma,\tau\in\mathfrak S_n}
   \operatorname{sgn}(\sigma)\operatorname{sgn}(\tau)
       \prod_{j=1}^n\langle u_{\sigma(j)},v_{\tau(j)}\rangle.
\]
For fixed $\sigma$, relabel $i=\sigma(j)$ and put
$\rho=\tau\circ\sigma^{-1}$. The sign product is
$\operatorname{sgn}(\rho)$ and the inner sum is the determinant
of the displayed mixed Gram matrix. There are $n!$ choices of
$\sigma$. Taking $v=u$ proves the norm identity.
\hfill$\square$

For independent $u_i$, divide by $\sqrt{n!\det G_u}$ to obtain
a unit vector. In particular the coordinate determinant divided
by $\sqrt{n!}$ has norm one. BMM, Proposition 5.19(1), p.49,
with $q=1$ and $b=n$, identifies the coefficient of the
holomorphic wedge cocycle with precisely this determinant shape.
The other active Fock coordinates and all remaining compact
Gaussians may be normalized to norm one. The unitary Bargmann
transform preserves the resulting norm. Under the standard
doubling map,
\[
 \delta(f_\infty\otimes\bar f_\infty)(0)=\|f_\infty\|^2=1.
\]
This is a chosen complete normalization, not an assertion that
an unmodified source convention already has the number one.
For the actual mixed endpoint vector retain its value
$a=\Phi_\infty(0)$; a certified continuity estimate
$|a-1|\leq\eta<1$ gives $|a|\geq1-\eta$.

\paragraph{Dimensions in the doubled model.}
For the ball of dimension $n$, $\dim_E V=n+1$. Doubling the
rank-$n$ skew-Hermitian source gives a split source of rank $2n$.
The Siegel unipotent is the additive group
$N\simeq\operatorname{Herm}_n(E)$ over $F$. A Hermitian matrix has
$n$ diagonal $F$-entries and $n(n-1)/2$ off-diagonal $E$-entries;
therefore
\[
 d_N:=\dim_\Q N=n^2[F:\Q].
\]
The Schr\"odinger position space is $V^n$ over $E$, of dimension
$n(n+1)$, and hence its real dimension at all archimedean places is
\[
 D=2n(n+1)[F:\Q].
\]
These are different dimensions: $d_N$ controls the unipotent grid,
whereas $D$ controls the position-lattice Gaussian tail.
In the grid statement below, $N$ is viewed by restriction of
scalars as a $\Q$-vector group; thus $N(\Q)$ denotes its
$F$-rational Hermitian matrices and $N(\A_f)$ their finite
$F$-adelic counterparts.

\paragraph{The explicit unipotent grid.}
Use probability measures on the compact quotients, fix
$\Phi_f(0)=1$, and fix the splitting for which the $G=\U(V)$
action in the doubled model is untwisted pullback. Let $K_N$ be
a compact open additive subgroup of $N(\A_f)$ fixing $\Phi_f$.
Set
\[
 \Lambda_N=N(\Q)\cap K_N,
 \qquad B:\R^{d_N}\longrightarrow N(\R),
\]
where the columns of $B$ are a $\Z$-basis of this full lattice.
Suppose the function
$F(b)=T_\Phi(n(b),1_f)$ is $L$-Lipschitz on
$N(\R)/\Lambda_N$. For $N_0\geq1$ define the explicit grid
\[
 \Gamma_{N_0}=
 \left\{\sum_{j=1}^{d_N}\frac{a_j}{N_0}\lambda_j:
                         0\leq a_j<N_0\right\}.
\]
Then one of its $N_0^{d_N}$ elements satisfies
\begin{equation}\label{Q10-grid}
 |T_\Phi(1,n_f(-\gamma))|
    \geq |a|-\frac{L\|B\|\sqrt{d_N}}{2N_0}.
\end{equation}

\paragraph{Proof.}
Additive strong approximation gives
$N(\A_f)=N(\Q)+K_N$ and identifies the quotient by $K_N$
with $N(\R)/\Lambda_N$. The quotient measure is its probability
Haar measure. The general rank-$r$ statement A1, applied with
$r=n$, gives its mean $a=\Phi_\infty(0)$: Siegel-unipotent
averaging annihilates every nonzero Gram matrix, and anisotropy
of $V$ leaves only the zero tuple. Hence
$\max |F|\geq|a|$.
Round the coordinates of a maximum point modulo $\Z^{d_N}$ to
the nearest grid point. Their Euclidean error is at most
$\sqrt{d_N}/(2N_0)$, so their images differ on the torus by at
most $\|B\|\sqrt{d_N}/(2N_0)$. The Lipschitz and reverse-triangle
inequalities give the bound for $|F(\gamma)|$.
Rational theta automorphy gives
$F(\gamma)=T_\Phi(1,n_f(-\gamma))$, proving
\eqref{Q10-grid}. The finite Weil action is unitary and commutes
with $G(\A_f)$, so the translate preserves both the finite
Schwartz norm and every compact target level that fixed the total
input. A factorization into individual inputs can require a
smaller common level. \hfill$\square$

If $|a|\geq1/2$, the explicit choice
\[
 N_0\geq\max\{1,\lceil2L\|B\|\sqrt{d_N}\rceil\}
\]
makes the error at most $1/4$, so the selected total integral
has absolute value at least $1/4$. Suppose the exact finite
tensor, source-character and connected-component expansion is
\[
 I=\sum_\lambda b_\lambda J_\lambda,\qquad
 S_{\rm ext}=\sum_\lambda|b_\lambda|>0,
\]
where the geometric normalization is included in $b_\lambda$.
Then at least one actual $2n$-factor integrated period satisfies
\begin{equation}\label{Q10-extract}
 |J_\lambda|\geq\frac{1}{4S_{\rm ext}}.
\end{equation}
This is the triangle inequality. If the tensor decomposition has
coefficients $c_a$ and its $2n$ torus factors have at most
$N_{i,a}$ permitted characters, a valid cost before the geometric
conversion is
$\sum_a|c_a|\prod_{i=1}^{2n}N_{i,a}$.
Zero character components may be retained in the count. The
certificate permits the displayed finite input translation; it
is not a lower bound for every preassigned unaveraged period.

\paragraph{A fully specified Lipschitz bound.}
Here is the all-rank version of the sufficient envelope in Q7.
Fix a compact representative set for the anisotropic target
quotient and a position lattice of minimum norm at least $\rho$,
containing the rational support for all its finite components.
Suppose, in the chosen Euclidean coordinates,
\[
 |\omega(g)\Phi(x)|
      \leq C(1+\|x\|)^M e^{-\lambda\|x\|^2}
\]
on that set, with $C,\lambda>0$. Write the unipotent phase as
$\exp(2\pi i\langle b,Q(x)\rangle)$ and suppose
$\|Q(x)\|\leq c_Q\|x\|^2$. Then one may take
\begin{equation}\label{Q10-L}
 L=2\pi c_QC\sum_{j\geq1}(2j+3)^D
       ((j+1)\rho)^2(1+(j+1)\rho)^M
                   e^{-\lambda j^2\rho^2}.
\end{equation}
To check this, differentiate the phase in a unit direction.
Its derivative has absolute value at most
$2\pi c_Q\|x\|^2$. The zero term contributes zero.
The shell $j\rho\leq\|x\|<(j+1)\rho$ contains at most
$(2j+3)^D$ points by the disjoint-ball argument. The displayed
bound follows by summing absolute values and integrating against
probability measure. Gaussian decay makes the differentiated
series uniformly convergent and justifies the operations.
The positive terms have an explicit polynomial factor times
$e^{-\lambda\rho^2j^2}$; after any index for which their
successive ratios are bounded by a number $q<1$, their tail is
at most its first term divided by $1-q$. This supplies a finite
certified evaluation.

For the determinant polynomial and its doubled vector one may
use $M=2n$ after the Gaussian model and its linear transformations
have been fixed: partial Fourier transform and metaplectic
changes of a polynomial Gaussian preserve the bound on its
polynomial degree, though they change the Gaussian and its
coefficients. In coordinates $Q_a(x)=x^tH_ax$ one can take
$c_Q=(\sum_a\|H_a\|^2)^{1/2}$. The lattice, envelope constants,
and geometric conversion must actually be supplied for numerical
use; no unrecorded effective reduction-theory bound is assumed.

\paragraph{A pointwise certificate for the top-form norm.}
If an injective holomorphic chart contains $B_R\subset\C^n$ and
$\xi=f(z)\,dz_1\wedge\cdots\wedge dz_n$ there, suppose
$|f(0)|\geq b>0$ and the real derivative of $f$ has norm at most
$M_1$ on the ball. For $0<r<R$ with $M_1r\leq b/2$,
\begin{equation}\label{Q10-topnorm}
 i^{n^2}\int_X\xi\wedge\bar\xi
       \geq \frac{2^{n-2}\pi^n}{n!}\,b^2r^{2n}.
\end{equation}
Indeed the segment estimate gives $|f|\geq b/2$ on $B_r$.
The positive top-form density is $2^n|f|^2$ times real
Lebesgue measure, and
$\operatorname{vol}(B_r\subset\R^{2n})=\pi^nr^{2n}/n!$.
Integration proves the formula. A determinant of first-tensor
values can be certified from individual errors $\varepsilon_i$
using the multilinear estimate
\[
 \left\|\bigwedge_i u_i-\bigwedge_i u_i^0\right\|
 \leq\prod_i(\|u_i^0\|+\varepsilon_i)-\prod_i\|u_i^0\|.
\]
Expand the wedge by subsets and apply the norm bound for each
term to prove this inequality. Together with A6's lattice tail,
it gives explicit input data for \eqref{Q10-topnorm}.

\paragraph{The definite subgroup and its local constant.}
Choose an $E$-rational line $L$ negative at the active place.
Then $H=\U(L^\perp)$ is totally definite of rank $n$. For a
canonical $n$-form, its right-regular coefficient at that line
has compact character
\[
 \kappa(\operatorname{diag}(A,d))=\det(A)d^{-n}.
\]
This follows because the tangent representation is
$\Hom(L,Q)$, whose derivative is $A/d$ and whose determinant
is $\det(A)d^{-n}$. It is trivial on the scalar center.
Suppose a unit vector $w$ in a unitary representation has that
compact type. With $\pi(A)=\det A$ on $H_\infty=\U(n)$ and
probability Haar,
\[
 \int_{\U(n)}
   \langle w,\sigma(\operatorname{diag}(A,1))w\rangle
                      \pi(A)^{-1}\,dA=1.
\]
The matrix coefficient is $\det A$ in the convention linear
in the second variable, so the integrand is one. The rank-zero
unipotent is trivial. This direct local computation is valid
for every $n$; the particular archive proposition imposing
even $k=n$ is invoked only when its additional hypotheses hold.
No global norm is identified with this local number.

\paragraph{Increasing the rank does not enlarge the canonical gap.}
For the compact sign pattern $(+,\ldots,+,-)$, the half sums are
\[
 \rho_c=\left(\frac{n-1}{2},\frac{n-3}{2},\ldots,
                         \frac{1-n}{2};0\right),\qquad
 \rho_n=\left(\frac12,\ldots,\frac12;-\frac n2\right).
\]
The right-regular canonical lowest compact weight is
$(1,\ldots,1;-n)$. Solving the archive's lowest-weight relation
$\lambda+\rho_n-\rho_c=(1,\ldots,1;-n)$ gives
\[
 \lambda=(1,\ldots,1;-n)-\rho_n+\rho_c
       =\left(\frac n2,\frac n2-1,\ldots,
                          1-\frac n2;-\frac n2\right).
\]
The compact determinant representation has tuple
$u_i=(n+3-2i)/2$, $1\leq i\leq n$.
Taking $v$ to be the first $n-1$ entries of $\lambda$ leaves
the archive's two nodes
\[
 b=\left(\frac n2-1,\frac n2\right).
\]
Their gap is exactly one in every rank.
The interlacing follows from
$u_i=\lambda_i+1/2>\lambda_i>u_{i+1}$ and the final inequality
$u_n>-\!b_0$; all cross differences are half-integral.
Thus increasing the dimension does not repair the canonical
\texttt{NodeGap} condition, whose required gap is at least
three. Invoking a general discrete-series theorem also requires
the representation identification demanded by that theorem;
the displayed direct compact integral does not require it.

For a pluricanonical power $K_X^q$, $q\geq1$, the compact weight
is $(q,\ldots,q;-nq)$. The same computation gives
\[
 \lambda_i=q+\frac n2-i\quad(1\leq i\leq n),\qquad
 \lambda_{n+1}=-nq+\frac n2,
 \qquad
 b=\left(\frac n2-q,nq-\frac n2\right).
\]
Its gap is $(n+1)q-n$. For $n=2$ this is $3q-2$, as in B9.
It enlarges the gap by changing the line bundle and its vector;
it does not identify the resulting section with the original
ordinary degree-$n$ cohomology class.
